\subsection{Further topics on random variables}


\subsubsection{Functions of a discrete random variable}

\textsc{PMF of a general function}

Let \(X\) be a discrete random variable with PMF \(p_X(x)\), and let \(Y\) be defined as a function of \(X\), say
\[
Y = g(X).
\]
We wish to find the PMF \(p_Y(y)\) of \(Y\).  
For any value \(y\) in the range of \(Y\), the event \(\{Y = y\}\) is the same as the event that \(X\) takes any value \(x\) such that \(g(x) = y\).  
Therefore,
\[
p_Y(y)
=
\mathbb{P}(Y = y)
=
\mathbb{P}\big(g(X) = y\big)
=
\sum_{x :\, g(x) = y} \mathbb{P}(X = x)
=
\sum_{x :\, g(x) = y} p_X(x).
\]
In words, to obtain \(p_Y(y)\), we sum the probabilities of all values of \(X\) that are mapped to \(y\) by the function \(g\).

\textsc{Example: many-to-one mapping}

Suppose that \(X\) takes the values \(2,3,4,5\) with probabilities
\[
p_X(2) = 0.1,
\quad
p_X(3) = 0.2,
\quad
p_X(4) = 0.3,
\quad
p_X(5) = 0.4,
\]
and define \(Y = g(X)\) by specifying \(g(2), g(3), g(4), g(5)\).  
If, for instance, \(g(4) = 4\) and \(g(5) = 4\), then the probability that \(Y\) takes the value \(4\) is
\[
p_Y(4)
=
\mathbb{P}(Y = 4)
=
\mathbb{P}(X = 4) + \mathbb{P}(X = 5)
=
p_X(4) + p_X(5)
=
0.3 + 0.4
=
0.7.
\]
This illustrates the general rule that whenever several values of \(X\) are mapped to the same value \(y\), the corresponding probability masses add up in the PMF of \(Y\).

\textsc{Linear functions: stretching and shifting}

Consider now the special case of a linear function
\[
Y = aX + b,
\]
where \(a\) and \(b\) are constants and \(a \ne 0\).  
For any value \(y\), the event \(\{Y = y\}\) is the same as the event that \(X\) takes the unique value
\[
x = \frac{y - b}{a}
\]
for which \(a x + b = y\), provided that this \(x\) lies in the support of \(X\).  
Thus,
\[
p_Y(y)
=
\mathbb{P}(Y = y)
=
\mathbb{P}\Big(X = \frac{y - b}{a}\Big)
=
p_X\Big(\frac{y - b}{a}\Big),
\]
whenever \((y - b)/a\) is a possible value of \(X\), and \(p_Y(y) = 0\) otherwise.

\textsc{Example: scaling}

Let \(X\) be a discrete random variable taking the values \(-1,1,2\) with
\[
p_X(-1) = \frac{2}{6},
\quad
p_X(1) = \frac{1}{6},
\quad
p_X(2) = \frac{3}{6}.
\]
Define
\[
Z = 2X.
\]
Then the possible values of \(Z\) are
\[
-2 \text{ when } X = -1,
\quad
2 \text{ when } X = 1,
\quad
4 \text{ when } X = 2.
\]
Using the general rule,
\[
p_Z(-2)
=
\mathbb{P}(Z = -2)
=
\mathbb{P}(X = -1)
=
\frac{2}{6},
\]
\[
p_Z(2)
=
\mathbb{P}(Z = 2)
=
\mathbb{P}(X = 1)
=
\frac{1}{6},
\]
\[
p_Z(4)
=
\mathbb{P}(Z = 4)
=
\mathbb{P}(X = 2)
=
\frac{3}{6}.
\]
If we plot the PMF of \(Z\) as a function of \(z\), it has exactly the same shape as the PMF of \(X\), except that the mass points have been stretched horizontally by a factor of \(2\).  
No probabilities change, only the locations of the mass points on the horizontal axis.

\textsc{Example: scaling and shifting}

Using the previous example, define
\[
Y = 2X + 3 = Z + 3.
\]
The possible values of \(Y\) are obtained by applying this transformation to the values of \(X\):
\[
X = -1 \Rightarrow Y = 1,
\quad
X = 1 \Rightarrow Y = 5,
\quad
X = 2 \Rightarrow Y = 7.
\]
The corresponding PMF values are
\[
p_Y(1)
=
\mathbb{P}(Y = 1)
=
\mathbb{P}(X = -1)
=
\frac{2}{6},
\]
\[
p_Y(5)
=
\mathbb{P}(Y = 5)
=
\mathbb{P}(X = 1)
=
\frac{1}{6},
\]
\[
p_Y(7)
=
\mathbb{P}(Y = 7)
=
\mathbb{P}(X = 2)
=
\frac{3}{6}.
\]
Thus the PMF of \(Y\) has the same shape as the PMF of \(Z\), but all mass points have been shifted to the right by \(3\).  
In general, for \(Y = aX + b\), the factor \(a\) stretches or compresses the horizontal axis, and the constant \(b\) shifts the PMF left or right, while the probabilities attached to corresponding mass points remain unchanged.

\textsc{General formula for a linear function}

For a discrete random variable \(X\) with PMF \(p_X\) and a linear function
\[
Y = aX + b,
\quad a \ne 0,
\]
the PMF of \(Y\) is given by
\[
p_Y(y)
=
p_X\Big(\frac{y - b}{a}\Big),
\]
for those values of \(y\) such that \((y - b)/a\) is in the support of \(X\), and \(p_Y(y) = 0\) for all other values of \(y\).  
This formula captures algebraically the graphical operations of stretching the PMF of \(X\) by the factor \(a\) and then shifting it horizontally by \(b\).



\subsubsection{A linear function of a continuous random variable}

\textsc{Heuristic picture}

Let \(X\) be a continuous random variable with PDF \(f_X\).  
Consider a linear transformation
\[
Y = aX + b,
\]
where \(a\) and \(b\) are constants and \(a \neq 0\).  
Geometrically, this transformation takes the horizontal axis, stretches it by a factor of \(a\), and then shifts it by \(b\).  
If we think of the PDF as a curve whose area under it must stay equal to \(1\), then stretching the horizontal axis by a factor \(|a|\) forces us to compress the vertical scale by a factor \(|a|\) so that the total area remains \(1\).

\textsc{Intuitive example}

Suppose that \(X\) has a piecewise constant PDF supported on the interval \([-1,1]\).  
Define
\[
Z = 2X.
\]
Then \(Z\) takes values in \([-2,2]\).  
The subinterval \([-1,0]\) of \(X\) corresponds to \([-2,0]\) for \(Z\), and if \(f_X\) is constant on \([-1,0]\), then the values of \(X\) in this interval are equally likely in the sense of density, so the corresponding values of \(Z\) in \([-2,0]\) should also have a constant density.  
Similarly, the subinterval \([0,1]\) of \(X\) corresponds to \([0,2]\) of \(Z\), and if \(f_X\) is constant on \([0,1]\), then \(f_Z\) is constant on \([0,2]\).

Assume, for instance, that the probability that \(X\) is positive equals \(2/3\), which is the area of a rectangle of width \(1\) and height \(2/3\).  
When we pass to \(Z\), the corresponding probability that \(Z\) is positive must still be \(2/3\), but now this probability is the area of a rectangle of width \(2\), so its height must be \(1/3\).  
The same reasoning applied to the negative side shows that if the area for negative \(X\) was \(1/3\), then the rectangle for negative \(Z\) now has width \(2\), so its height must be \(1/6\).  
Thus, the shape of the PDF is stretched horizontally by a factor of \(2\), but the heights are scaled down by a factor of \(2\) to preserve total area \(1\).

Now add a constant and define
\[
Y = Z + 3 = 2X + 3.
\]
Then the interval \([-2,0]\) for \(Z\) becomes \([1,3]\) for \(Y\), and \([0,2]\) becomes \([3,5]\).  
The areas of the corresponding rectangles must remain the same, so their heights are unchanged.  
In other words, adding a constant simply shifts the PDF horizontally without altering its shape or vertical scale.

\textsc{CDF-based derivation for \(a>0\)}

We now derive a general formula.  
Assume first that \(a>0\).  
The CDF of \(Y\) is
\[
F_Y(y)
=
\mathbb{P}(Y \le y)
=
\mathbb{P}(aX + b \le y).
\]
Solving the inequality for \(X\) gives
\[
aX + b \le y
\quad\Longleftrightarrow\quad
X \le \frac{y - b}{a}.
\]
Therefore,
\[
F_Y(y)
=
\mathbb{P}\Big(X \le \frac{y - b}{a}\Big)
=
F_X\Big(\frac{y - b}{a}\Big),
\]
where \(F_X\) is the CDF of \(X\).  
To obtain the PDF, we differentiate both sides with respect to \(y\).  
Using the chain rule,
\[
f_Y(y)
=
\frac{d}{dy} F_Y(y)
=
\frac{d}{dy} F_X\Big(\frac{y - b}{a}\Big)
=
f_X\Big(\frac{y - b}{a}\Big)\cdot\frac{d}{dy}\Big(\frac{y - b}{a}\Big)
=
\frac{1}{a}\,f_X\Big(\frac{y - b}{a}\Big),
\]
for all \(y\) such that \((y-b)/a\) lies in the support of \(X\), and \(f_Y(y)=0\) otherwise.

\textsc{CDF-based derivation for \(a<0\)}

Assume now that \(a<0\).  
We again start from
\[
F_Y(y)
=
\mathbb{P}(Y \le y)
=
\mathbb{P}(aX + b \le y).
\]
Solving the inequality for \(X\), the sign of the inequality reverses when we divide by the negative constant \(a\):
\[
aX + b \le y
\quad\Longleftrightarrow\quad
X \ge \frac{y - b}{a}.
\]
Thus,
\[
F_Y(y)
=
\mathbb{P}\Big(X \ge \frac{y - b}{a}\Big)
=
1 - \mathbb{P}\Big(X < \frac{y - b}{a}\Big).
\]
Since \(X\) is continuous, \(\mathbb{P}(X < c) = \mathbb{P}(X \le c)\) for any constant \(c\), so
\[
F_Y(y)
=
1 - F_X\Big(\frac{y - b}{a}\Big).
\]
Differentiating both sides with respect to \(y\) and again using the chain rule gives
\[
f_Y(y)
=
\frac{d}{dy} F_Y(y)
=
- f_X\Big(\frac{y - b}{a}\Big)\cdot\frac{d}{dy}\Big(\frac{y - b}{a}\Big)
=
- f_X\Big(\frac{y - b}{a}\Big)\cdot\frac{1}{a}
=
\frac{1}{|a|}\,f_X\Big(\frac{y - b}{a}\Big),
\]
because when \(a<0\) we have \(-1/a = 1/|a|\).

\textsc{Unified formula and interpretation}

Combining the two cases \(a>0\) and \(a<0\), we obtain the general formula for any \(a\neq 0\):
\[
f_Y(y)
=
\frac{1}{|a|}\,
f_X\Big(\frac{y - b}{a}\Big),
\]
for all \(y\) such that \((y-b)/a\) lies in the support of \(X\), and \(f_Y(y)=0\) otherwise.  
This formula encodes three geometric effects of the transformation \(Y = aX + b\):
\begin{itemize}
\item The argument \((y-b)/a\) inside \(f_X\) corresponds to a horizontal stretching or compression by the factor \(a\).  
\item The term \(b\) corresponds to a horizontal shift of the PDF by \(b\).  
\item The factor \(1/|a|\) corresponds to a vertical scaling that ensures the area under \(f_Y\) is equal to \(1\).  
\end{itemize}

\textsc{Comparison with the discrete case}

In the discrete case, for \(Y = aX + b\) we had
\[
p_Y(y)
=
p_X\Big(\frac{y - b}{a}\Big),
\]
without any extra scaling factor.  
This reflects the fact that in the discrete setting we preserve the individual point probabilities exactly, and there is no notion of stretching intervals and adjusting heights to keep total probability equal to \(1\).  
In the continuous setting, probabilities are represented by areas under the PDF, so when we stretch the horizontal axis by \(|a|\), we must divide the height of the curve by \(|a|\) to keep the total area unchanged.



\subsubsection{A linear function of a normal random variable}

\textsc{Setup}

Let \(X\) be a normal random variable with mean \(\mu\) and variance \(\sigma^{2}\), so
\[
X \sim \mathcal{N}(\mu,\sigma^{2}),
\]
and its PDF is
\[
f_{X}(x)
=
\frac{1}{\sqrt{2\pi}\,\sigma}
\exp\!\Bigg(
-\frac{1}{2}
\Bigg(
\frac{x-\mu}{\sigma}
\Bigg)^{2}
\Bigg).
\]
Consider the linear transformation
\[
Y = aX + b,
\]
where \(a\) and \(b\) are constants with \(a \neq 0\).  

\textsc{Using the general transformation formula}

From the general result for linear functions of a continuous random variable, the PDF of \(Y\) is given by
\[
f_{Y}(y)
=
\frac{1}{|a|}\,
f_{X}\Big(\frac{y-b}{a}\Big).
\]
Substituting the normal PDF of \(X\) into this expression gives
\[
f_{Y}(y)
=
\frac{1}{|a|}
\cdot
\frac{1}{\sqrt{2\pi}\,\sigma}
\exp\!\Bigg(
-\frac{1}{2}
\Bigg(
\frac{\frac{y-b}{a}-\mu}{\sigma}
\Bigg)^{2}
\Bigg).
\]
We now simplify the expression inside the exponential and the prefactor.  
First rewrite
\[
\frac{\frac{y-b}{a}-\mu}{\sigma}
=
\frac{y-b-a\mu}{a\sigma},
\]
so the exponent becomes
\[
-\frac{1}{2}
\Bigg(
\frac{y-b-a\mu}{a\sigma}
\Bigg)^{2}
=
-\frac{1}{2}
\Bigg(
\frac{y-(a\mu+b)}{|a|\sigma}
\Bigg)^{2},
\]
because the square absorbs the sign of \(a\).  

For the prefactor, we have
\[
\frac{1}{|a|}
\cdot
\frac{1}{\sqrt{2\pi}\,\sigma}
=
\frac{1}{\sqrt{2\pi}\,|a|\sigma}
=
\frac{1}{\sqrt{2\pi}\,\sqrt{a^{2}\sigma^{2}}}.
\]
Putting these together,
\[
f_{Y}(y)
=
\frac{1}{\sqrt{2\pi}\,\sqrt{a^{2}\sigma^{2}}}
\exp\!\Bigg(
-\frac{1}{2}
\Bigg(
\frac{y-(a\mu+b)}{\sqrt{a^{2}\sigma^{2}}}
\Bigg)^{2}
\Bigg).
\]

\textsc{Identification of the resulting distribution}

The last expression has exactly the standard normal PDF form in the variable \(y\), with mean
\[
\mu_{Y} = a\mu + b
\]
and variance
\[
\sigma_{Y}^{2} = a^{2}\sigma^{2}.
\]
Therefore,
\[
Y = aX + b
\sim
\mathcal{N}(a\mu + b,\; a^{2}\sigma^{2}).
\]
The fact that the mean and variance transform as \(E[Y] = a\mu + b\) and \(\operatorname{Var}(Y) = a^{2}\sigma^{2}\) agrees with the general formulas for linear functions of any random variable.  
The special feature here is that the entire distributional shape is preserved: the bell-shaped normal PDF of \(X\) is stretched and shifted, but remains bell shaped, so the resulting random variable \(Y\) is again normal.  


\subsubsection{The PDF of a general function of a continuous r.v.}

\textsc{Two-step method}

Let \(X\) be a continuous random variable with known PDF \(f_X\) and CDF \(F_X\), and let
\[
Y = g(X)
\]
for some (possibly nonlinear) function \(g\).  
To find the PDF of \(Y\), we use a two-step procedure:
\begin{enumerate}
\item First find the CDF of \(Y\),
\[
F_Y(y)
=
\mathbb{P}(Y \le y)
=
\mathbb{P}\big(g(X) \le y\big),
\]
by rewriting the event \(\{g(X) \le y\}\) in terms of inequalities involving \(X\) and then computing that probability using the known distribution of \(X\).  
\item Then obtain the PDF by differentiation,
\[
f_Y(y)
=
\frac{d}{dy}F_Y(y),
\]
for those values of \(y\) in the range of \(Y\).  
\end{enumerate}
Most of the work lies in the first step, where we translate the condition \(g(X)\le y\) into a suitable interval (or union of intervals) for \(X\) and then evaluate the corresponding probability.

\textsc{Example: \(Y = X^{3}\) with \(X \sim \text{Uniform}(0,2)\)}

Suppose that \(X\) is uniform on \([0,2]\), so its PDF is
\[
f_X(x)
=
\begin{cases}
\frac{1}{2}, & 0 \le x \le 2,\\[4pt]
0, & \text{otherwise}.
\end{cases}
\]
Define
\[
Y = X^{3}.
\]
Since \(X\) takes values in \([0,2]\), the range of \(Y\) is \([0,8]\).  
For \(y<0\) we have \(F_Y(y)=0\), and for \(y>8\) we have \(F_Y(y)=1\), so we focus on \(0 \le y \le 8\).

For such \(y\),
\[
F_Y(y)
=
\mathbb{P}(Y \le y)
=
\mathbb{P}(X^{3} \le y).
\]
On \([0,2]\), the function \(x \mapsto x^{3}\) is increasing, so the inequality \(X^{3} \le y\) is equivalent to
\[
X \le y^{1/3}.
\]
Thus,
\[
F_Y(y)
=
\mathbb{P}\big(X \le y^{1/3}\big)
=
\int_{0}^{y^{1/3}} \frac{1}{2}\,dx
=
\frac{1}{2}\,y^{1/3},
\quad 0 \le y \le 8.
\]
This completes the first step.  
For the second step, we differentiate:
\[
f_Y(y)
=
\frac{d}{dy}F_Y(y)
=
\frac{d}{dy}\Big(\frac{1}{2}y^{1/3}\Big)
=
\frac{1}{2}\cdot\frac{1}{3}y^{-2/3}
=
\frac{1}{6}\,y^{-2/3},
\quad 0 < y \le 8.
\]
For \(y\) outside \([0,8]\), we have \(f_Y(y)=0\).  
So the PDF of \(Y\) is
\[
f_Y(y)
=
\begin{cases}
\frac{1}{6}\,y^{-2/3}, & 0<y\le 8,\\[4pt]
0, & \text{otherwise}.
\end{cases}
\]
This PDF is not constant: it becomes larger as \(y\) approaches \(0\), and in fact \(f_Y(y)\) tends to infinity as \(y \downarrow 0\), while still integrating to \(1\) over \((0,8]\).

\textsc{Example: \(Y = 10/X\) with \(X \sim \text{Uniform}(5,10)\)}

Now let \(X\) be uniform on \([5,10]\), so that
\[
f_X(x)
=
\begin{cases}
\frac{1}{5}, & 5 \le x \le 10,\\[4pt]
0, & \text{otherwise}.
\end{cases}
\]
Define the random variable
\[
Y = \frac{10}{X}.
\]
Since \(X \in [5,10]\), the smallest value of \(Y\) occurs when \(X\) is largest, \(X=10\), giving \(Y=1\), and the largest value of \(Y\) occurs when \(X\) is smallest, \(X=5\), giving \(Y=2\).  
Thus, the range of \(Y\) is \([1,2]\), and \(f_Y(y)=0\) outside this interval.

For \(1 \le y \le 2\),
\[
F_Y(y)
=
\mathbb{P}(Y \le y)
=
\mathbb{P}\Big(\frac{10}{X} \le y\Big).
\]
Because \(X>0\), the inequality \(\frac{10}{X} \le y\) is equivalent to
\[
10 \le yX
\quad\Longleftrightarrow\quad
X \ge \frac{10}{y}.
\]
So
\[
F_Y(y)
=
\mathbb{P}\Big(X \ge \frac{10}{y}\Big).
\]
For \(1 \le y \le 2\), the point \(10/y\) lies in \([5,10]\), so we can express this probability as the area under the uniform density from \(10/y\) to \(10\):
\[
F_Y(y)
=
\int_{10/y}^{10} \frac{1}{5}\,dx
=
\frac{1}{5}\Big(10 - \frac{10}{y}\Big)
=
\frac{10}{5}\Big(1 - \frac{1}{y}\Big)
=
2\Big(1 - \frac{1}{y}\Big),
\quad 1 \le y \le 2.
\]
Outside this interval, \(F_Y(y)\) is \(0\) for \(y<1\) and \(1\) for \(y>2\).

To obtain the PDF, we differentiate on \((1,2)\):
\[
f_Y(y)
=
\frac{d}{dy}F_Y(y)
=
\frac{d}{dy}\Big(2 - \frac{2}{y}\Big)
=
0 - 2\cdot\Big(-\frac{1}{y^{2}}\Big)
=
\frac{2}{y^{2}},
\quad 1<y<2.
\]
Thus,
\[
f_Y(y)
=
\begin{cases}
\displaystyle \frac{2}{y^{2}}, & 1<y<2,\\[4pt]
0, & \text{otherwise}.
\end{cases}
\]
The PDF starts at a value of \(2\) when \(y\) is close to \(1\) and decreases as \(y\) increases, reflecting that shorter times (corresponding to higher speeds) are more likely than longer times in this setup.

\textsc{General strategy}

These examples illustrate a general and systematic approach to finding the PDF of a function \(Y = g(X)\) of a continuous random variable \(X\):
\begin{itemize}
\item Identify the range of \(Y\), so that you know where \(f_Y(y)\) can be nonzero.  
\item For \(y\) in this range, express the event \(\{Y \le y\}\) as an event involving \(X\), typically of the form \(\{X \le h(y)\}\), \(\{X \ge h(y)\}\), or a union of such intervals.  
\item Compute \(F_Y(y)\) by integrating the known PDF \(f_X\) over the appropriate interval(s) for \(X\).  
\item Differentiate \(F_Y(y)\) with respect to \(y\) to obtain \(f_Y(y)\).  
\end{itemize}
In more complicated situations where \(g\) is not monotone on the entire support of \(X\), one splits the support of \(X\) into regions on which \(g\) is monotone, computes the contribution from each region, and then adds them up to obtain the CDF or PDF of \(Y\).


\subsubsection{The monotonic case}

\textsc{Setup and inverse function}

Let \(X\) be a continuous random variable with PDF \(f_X\), and let
\[
Y = g(X),
\]
where \(g\) is a strictly monotonic and differentiable function on the range of interest of \(X\).  
Because \(g\) is strictly monotonic, every value \(y\) in the range of \(Y\) comes from a unique value \(x\), so \(g\) has an inverse function \(h\) satisfying
\[
x = h(y)
\quad\Longleftrightarrow\quad
y = g(x).
\]
The function \(h\) maps values of \(y\) back to the corresponding values of \(x\).

\textsc{Strictly increasing case}

Assume first that \(g\) is strictly increasing.  
Fix some \(y\) in the range of \(Y\).  
The CDF of \(Y\) is
\[
F_Y(y)
=
\mathbb{P}(Y \le y)
=
\mathbb{P}\big(g(X) \le y\big).
\]
Since \(g\) is increasing, the inequality \(g(X) \le y\) is equivalent to
\[
X \le h(y),
\]
because \(h(y)\) is the unique value of \(x\) that produces \(y\).  
Therefore,
\[
F_Y(y)
=
\mathbb{P}\big(X \le h(y)\big)
=
F_X\big(h(y)\big),
\]
where \(F_X\) is the CDF of \(X\).  
To obtain the PDF, we differentiate both sides with respect to \(y\) and apply the chain rule:
\[
f_Y(y)
=
\frac{d}{dy}F_Y(y)
=
\frac{d}{dy}F_X\big(h(y)\big)
=
f_X\big(h(y)\big)\,h'(y),
\]
for \(y\) in the range of \(Y\), and \(f_Y(y)=0\) otherwise.  
When \(g\) is increasing, the inverse \(h\) is also increasing, so \(h'(y)\ge 0\).

\textsc{Strictly decreasing case}

Assume now that \(g\) is strictly decreasing.  
Again fix \(y\) in the range of \(Y\), and start from
\[
F_Y(y)
=
\mathbb{P}(Y \le y)
=
\mathbb{P}\big(g(X) \le y\big).
\]
Because \(g\) is decreasing, the inequality \(g(X) \le y\) holds when \(X\) is larger than or equal to the value \(h(y)\) that produces \(y\), so
\[
F_Y(y)
=
\mathbb{P}\big(X \ge h(y)\big).
\]
This can be written in terms of the CDF of \(X\) as
\[
F_Y(y)
=
1 - \mathbb{P}\big(X < h(y)\big).
\]
Since \(X\) is continuous, \(\mathbb{P}(X < h(y)) = \mathbb{P}(X \le h(y)) = F_X\big(h(y)\big)\), and thus
\[
F_Y(y)
=
1 - F_X\big(h(y)\big).
\]
Differentiating both sides with respect to \(y\) gives
\[
f_Y(y)
=
\frac{d}{dy}F_Y(y)
=
- f_X\big(h(y)\big)\,h'(y).
\]
In the decreasing case, \(h\) is also decreasing, so \(h'(y) \le 0\).  
Thus the product \(-h'(y)\) is nonnegative, as a PDF must be.  

\textsc{Unified formula}

In the increasing case we have
\[
f_Y(y)
=
f_X\big(h(y)\big)\,h'(y),
\]
and in the decreasing case we have
\[
f_Y(y)
=
- f_X\big(h(y)\big)\,h'(y).
\]
These two expressions can be combined into a single formula using absolute values:
\[
f_Y(y)
=
f_X\big(h(y)\big)\,\big|h'(y)\big|,
\]
for \(y\) in the range of \(Y\), and \(f_Y(y)=0\) otherwise.  
This is the general transformation rule for the PDF of a strictly monotonic and differentiable function of a continuous random variable.

\textsc{Example: \(Y = X^{2}\) with \(X \sim \text{Uniform}(0,1)\)}

Let \(X\) be uniform on \([0,1]\), so
\[
f_X(x)
=
\begin{cases}
1, & 0 \le x \le 1,\\
0, & \text{otherwise},
\end{cases}
\]
and consider
\[
Y = g(X) = X^{2}.
\]
The function \(g(x)=x^{2}\) is not monotonic on the entire real line, but it is strictly increasing on \([0,1]\), which is the support of \(X\).  
Therefore we can apply the monotonic transformation formula on this interval.

First identify the inverse function.  
On \([0,1]\),
\[
y = x^{2}
\quad\Longleftrightarrow\quad
x = \sqrt{y},
\]
so
\[
h(y) = \sqrt{y}.
\]
Since \(X\in[0,1]\), \(Y = X^{2}\) also takes values in \([0,1]\), and outside this interval we must have \(f_Y(y)=0\).  
For \(0<y<1\),
\[
f_Y(y)
=
f_X\big(h(y)\big)\,\big|h'(y)\big|
=
f_X(\sqrt{y})\cdot\Big|\frac{d}{dy}\sqrt{y}\Big|.
\]
On the support of \(X\), we have \(f_X(\sqrt{y})=1\), and
\[
\frac{d}{dy}\sqrt{y}
=
\frac{1}{2\sqrt{y}}.
\]
Thus,
\[
f_Y(y)
=
1 \cdot \frac{1}{2\sqrt{y}}
=
\frac{1}{2\sqrt{y}},
\quad 0<y<1.
\]
Including the behavior outside \([0,1]\), the PDF of \(Y\) is
\[
f_Y(y)
=
\begin{cases}
\displaystyle \frac{1}{2\sqrt{y}}, & 0<y<1,\\[4pt]
0, & \text{otherwise}.
\end{cases}
\]
This result is obtained with much less work than going through the full two-step CDF method, because we only need to identify the inverse function \(h\) and its derivative.


\bigskip
\textsc{Matching small intervals}

Let \(Y = g(X)\), where \(g\) is a strictly monotonic and differentiable function, and let \(h\) be its inverse, so that
\[
x = h(y)
\quad\Longleftrightarrow\quad
y = g(x).
\]
Fix a particular value \(x\) and the corresponding value \(y = g(x)\).  
Consider a small interval around \(x\), say \([x, x + \Delta_1]\), and look at the corresponding interval for \(Y\), which is approximately \([y, y + \Delta_2]\), where \(\Delta_2\) is the small change in \(y\) that results from the change \(\Delta_1\) in \(x\).  
Because the mapping \(y = g(x)\) is one-to-one, the event
\[
\{X \in [x, x + \Delta_1]\}
\]
is the same as the event
\[
\{Y \in [y, y + \Delta_2]\},
\]
so these two events have the same probability:
\[
\mathbb{P}\big(y \le Y \le y + \Delta_2\big)
=
\mathbb{P}\big(x \le X \le x + \Delta_1\big).
\]

\textsc{Approximating probabilities with PDFs}

For small intervals, probabilities can be approximated using the PDFs.  
On the \(y\)-side,
\[
\mathbb{P}\big(y \le Y \le y + \Delta_2\big)
\approx
f_Y(y)\,\Delta_2,
\]
and on the \(x\)-side,
\[
\mathbb{P}\big(x \le X \le x + \Delta_1\big)
\approx
f_X(x)\,\Delta_1.
\]
Equating these two approximations gives
\[
f_Y(y)\,\Delta_2
\approx
f_X(x)\,\Delta_1.
\]
To turn this into a relation between \(f_Y\) and \(f_X\), we need to relate \(\Delta_1\) and \(\Delta_2\).

\textsc{Relating the interval lengths via the inverse}

The change \(\Delta_2\) in \(y\) is caused by the change \(\Delta_1\) in \(x\) through the function \(g\).  
Equivalently, we can look at how a small change \(\Delta_2\) in \(y\) produces a change \(\Delta_1\) in \(x\) through the inverse function \(h\).  
When \(y\) advances from \(y\) to \(y + \Delta_2\), the corresponding \(x\) advances from \(x\) to \(x + \Delta_1\), so
\[
\Delta_1
\approx
h'(y)\,\Delta_2,
\]
where \(h'(y)\) is the derivative (slope) of the inverse function at \(y\).  

Substituting this into the earlier approximate equality yields
\[
f_Y(y)\,\Delta_2
\approx
f_X(x)\,h'(y)\,\Delta_2.
\]
Canceling \(\Delta_2\) from both sides gives
\[
f_Y(y)
\approx
f_X(x)\,h'(y),
\]
and in the limit of vanishingly small intervals this becomes the exact relation
\[
f_Y(y)
=
f_X\big(h(y)\big)\,h'(y),
\]
where \(x = h(y)\) is the unique value of \(X\) that maps to \(y\).

\textsc{Absolute value and decreasing functions}

If \(g\) is decreasing instead of increasing, the same small-interval argument still applies.  
The difference is that in this case the inverse function \(h\) is also decreasing, so \(h'(y) \le 0\).  
The probability of a small interval in \(y\) must still be equal to the probability of the corresponding interval in \(x\), so the same algebra leads to
\[
f_Y(y)
=
f_X\big(h(y)\big)\,\big|h'(y)\big|.
\]
This absolute value ensures that the PDF \(f_Y\) is always nonnegative, regardless of whether \(g\) (and hence \(h\)) is increasing or decreasing.  
Thus, the intuitive picture is that small intervals on the \(x\)-axis are stretched or compressed when mapped to the \(y\)-axis, and the PDFs adjust inversely to this stretching so that the probabilities of corresponding small intervals remain the same.



\subsubsection{A function of multiple random variables}

\textsc{General methodology}

Let \(X\) and \(Y\) be jointly continuous random variables with known joint PDF \(f_{X,Y}(x,y)\).  
Suppose we define a new random variable
\[
Z = g(X,Y),
\]
for some function \(g\).  
To find the distribution of \(Z\), we follow the same two-step procedure used for functions of a single random variable:
\begin{enumerate}
\item First find the CDF of \(Z\),
\[
F_Z(z)
=
\mathbb{P}(Z \le z)
=
\mathbb{P}\big(g(X,Y) \le z\big)
=
\iint_{\{(x,y):\, g(x,y)\le z\}} f_{X,Y}(x,y)\,dx\,dy.
\]
\item Then obtain the PDF by differentiating,
\[
f_Z(z)
=
\frac{d}{dz}F_Z(z),
\]
for those values of \(z\) in the range of \(Z\).  
\end{enumerate}
In many examples, step one becomes a geometric problem: describe the region \(\{(x,y) : g(x,y)\le z\}\) and integrate \(f_{X,Y}\) over that region.

\textsc{Example: \(Z = Y/X\) with \((X,Y)\) uniform on the unit square}

Consider \(X\) and \(Y\) independent and each uniform on \([0,1]\).  
Their joint PDF is
\[
f_{X,Y}(x,y)
=
\begin{cases}
1, & 0 \le x \le 1,\; 0 \le y \le 1,\\[4pt]
0, & \text{otherwise}.
\end{cases}
\]
We define
\[
Z = \frac{Y}{X}
\]
and seek the CDF and PDF of \(Z\).

\textsc{CDF for negative and small positive \(z\)}

The CDF is
\[
F_Z(z)
=
\mathbb{P}(Z \le z)
=
\mathbb{P}\Big(\frac{Y}{X} \le z\Big).
\]
If \(z<0\), then \(\frac{Y}{X} \ge 0\) always, because \(X\) and \(Y\) are nonnegative, so
\[
F_Z(z) = 0, \quad z<0.
\]

Now consider \(0 \le z \le 1\).  
The inequality \(\frac{Y}{X} \le z\) is equivalent to
\[
Y \le zX.
\]
On the \((x,y)\)-plane, this is the region below the line \(y = zx\).  
When \(0 \le z \le 1\), this line intersects the top edge \(y=1\) at \(x = 1/z \ge 1\), so over the unit square it intersects only the right edge at \((1,z)\).  
Thus, the region \(\{(x,y): 0 \le x \le 1,\; 0 \le y \le z x\}\) is a right triangle with base \(1\) and height \(z\), so its area is
\[
\text{area} = \frac{1}{2}\cdot 1 \cdot z = \frac{z}{2}.
\]
Since the joint density is 1 on the unit square, this area equals the probability, and therefore
\[
F_Z(z) = \frac{z}{2}, \quad 0 \le z \le 1.
\]

\textsc{CDF for \(z>1\)}

Now consider \(z>1\).  
Again, the event \(\frac{Y}{X}\le z\) corresponds to the region below the line \(y = z x\).  
When \(z>1\), this line intersects the top edge \(y=1\) at \(x = 1/z < 1\), so it cuts off a small triangle near the origin where \(y > z x\).  
The region where \(Y/X \le z\) is then the whole unit square minus that triangle.

The small excluded triangle has base \(1/z\) on the \(x\)-axis and height \(1\) on the \(y\)-axis, so its area is
\[
\text{area}_{\text{triangle}}
=
\frac{1}{2}\cdot \frac{1}{z}\cdot 1
=
\frac{1}{2z}.
\]
Thus, the area of the region where \(Y/X \le z\) is
\[
1 - \frac{1}{2z},
\]
and consequently
\[
F_Z(z)
=
1 - \frac{1}{2z},
\quad z>1.
\]

Putting all cases together, the CDF of \(Z\) is
\[
F_Z(z)
=
\begin{cases}
0, & z<0,\\[4pt]
\displaystyle \frac{z}{2}, & 0 \le z \le 1,\\[6pt]
\displaystyle 1 - \frac{1}{2z}, & z>1.
\end{cases}
\]

\textsc{PDF of \(Z\)}

To obtain the PDF, we differentiate \(F_Z(z)\) piecewise.  

For \(z<0\), \(F_Z(z)\) is constant, so
\[
f_Z(z) = 0, \quad z<0.
\]

For \(0<z<1\),
\[
F_Z(z) = \frac{z}{2}
\quad\Longrightarrow\quad
f_Z(z)
=
\frac{d}{dz}F_Z(z)
=
\frac{1}{2}.
\]

For \(z>1\),
\[
F_Z(z)
=
1 - \frac{1}{2z}
\quad\Longrightarrow\quad
f_Z(z)
=
\frac{d}{dz}\Big(1 - \frac{1}{2z}\Big)
=
\frac{1}{2z^{2}}.
\]

At the transition points \(z=0\) and \(z=1\), the PDF can be defined by continuity on either side, but this does not affect probabilities because these are single points.  
Thus, the PDF of \(Z = Y/X\) is
\[
f_Z(z)
=
\begin{cases}
0, & z<0,\\[4pt]
\displaystyle \frac{1}{2}, & 0<z<1,\\[6pt]
\displaystyle \frac{1}{2z^{2}}, & z>1.
\end{cases}
\]

\textsc{Interpretation}

The distribution of \(Z = Y/X\) arises from the geometry of the uniform joint distribution on the unit square.  
For small positive \(z\), the allowed region where \(Y/X \le z\) is a triangle whose area grows linearly with \(z\), giving a constant PDF on \((0,1)\).  
For larger \(z\), the allowed region is the unit square minus a shrinking triangle near the origin, whose area decreases like \(1/z\), and the corresponding PDF decays like \(1/z^{2}\). 



\subsubsection{Convolution of PMFs for sums of independent random variables}

Let \(X\) and \(Y\) be independent discrete random variables with known PMFs \(p_X(x)\) and \(p_Y(y)\).  
We define a new random variable
\[
Z = X + Y
\]
and we want to find the PMF \(p_Z(z) = \mathbb{P}(Z = z)\).

\textsc{Basic idea}

The event \(\{Z = z\}\) means that the pair \((X,Y)\) takes values that add up to \(z\).  
There are many ways for this to happen; for example, if \(z=3\) then we could have \((X,Y) = (0,3)\), or \((1,2)\), or \((2,1)\), and so on.  
The probability that \(Z\) equals \(z\) is obtained by adding the probabilities of all such pairs.  
Because \(X\) and \(Y\) are independent, the probability of any particular pair \((x,y)\) is the product \(p_X(x)\,p_Y(y)\).  
This leads to a systematic formula.

\textsc{Convolution formula}

Fix an integer \(z\).  
For each possible value \(x\) of \(X\), the value of \(Y\) must be \(z-x\) in order for the sum to equal \(z\).  
Thus
\[
\mathbb{P}(Z = z)
=
\sum_{x} \mathbb{P}\big(X = x,\; Y = z-x\big).
\]
Independence gives
\[
\mathbb{P}\big(X = x,\; Y = z-x\big)
=
\mathbb{P}(X = x)\,\mathbb{P}(Y = z-x)
=
p_X(x)\,p_Y(z-x).
\]
Therefore the PMF of \(Z\) is
\[
p_Z(z)
=
\sum_{x} p_X(x)\,p_Y(z-x),
\]
where the sum runs over all values of \(x\) for which both terms in the product are defined and nonzero.  
This operation on PMFs is called \textit{convolution}.  
Given \(p_X\) and \(p_Y\), the convolution formula produces a new PMF \(p_Z\) corresponding to the sum \(Z = X+Y\).

\textsc{Example}

Suppose \(X\) takes values \(1\) and \(4\) with
\[
p_X(1) = \frac{1}{3},
\qquad
p_X(4) = \frac{2}{3},
\]
and \(Y\) takes values \(-2,-1,1\) with
\[
p_Y(-2) = \frac{3}{6},
\qquad
p_Y(-1) = \frac{2}{6},
\qquad
p_Y(1) = \frac{1}{6},
\]
and \(X\) and \(Y\) are independent.  
To compute \(p_Z(3)\) where \(Z = X+Y\), list all pairs \((x,y)\) with \(x+y=3\).  
These are
\[
(x,y) = (1,2) \textit{ (impossible here)}, \quad (4,-1).
\]
In this particular example the only feasible pair is \((4,-1)\) because \(Y\) never takes the value \(2\).  
Thus
\[
p_Z(3)
=
\mathbb{P}(X = 4,\; Y = -1)
=
p_X(4)\,p_Y(-1)
=
\frac{2}{3}\cdot\frac{2}{6}
=
\frac{4}{18}
=
\frac{2}{9}.
\]
If instead \(Y\) could also take the value \(2\) with some probability, say \(p_Y(2)\), then
\[
p_Z(3)
=
p_X(1)\,p_Y(2)
+
p_X(4)\,p_Y(-1)
=
\frac{1}{3}p_Y(2)
+
\frac{2}{3}\cdot\frac{2}{6}.
\]

\textsc{Graphical interpretation}

A useful graphical method to remember the convolution formula is to work with the PMFs drawn as stem plots on the number line.  
One can reflect the PMF of \(Y\) about the vertical axis to obtain a plot of \(p_Y(-y)\).  
Then shift this reflected plot by an amount \(z\) so that the point corresponding to \(y\) lines up with the point \(x = z-y\).  
For a given shift \(z\), the probability \(p_Z(z)\) is obtained by multiplying corresponding vertical stems that line up and summing all these products.  
Changing \(z\) corresponds to sliding the reflected PMF left or right and repeating the same multiplication and summation.  
This picture mirrors exactly the algebraic convolution sum \(p_Z(z) = \sum_x p_X(x)\,p_Y(z-x)\).



\subsubsection{Convolution for sums of independent continuous random variables}

Let \(X\) and \(Y\) be independent continuous random variables with PDFs \(f_X(x)\) and \(f_Y(y)\).  
Define
\[
Z = X + Y.
\]
We want to find the PDF \(f_Z(z)\) of the sum \(Z\).

\textsc{Heuristic derivation}

For a fixed value \(z\), the event \(\{Z = z\}\) corresponds to all pairs \((x,y)\) satisfying \(x + y = z\).  
In the discrete case, to obtain the PMF we add the probabilities of all pairs on that line and get a sum (the discrete convolution).  
In the continuous case, there are infinitely many pairs on the line \(x+y=z\), so the “sum over ways” becomes an integral and PMFs are replaced by PDFs.  
This leads to the continuous convolution formula
\[
f_Z(z)
=
\int_{-\infty}^{\infty} f_X(x)\,f_Y(z-x)\,dx,
\]
where we integrate over all possible values of \(x\) and for each such \(x\) the corresponding \(y\) is \(z-x\).

\textsc{Conditional-density derivation}

A more formal justification starts by conditioning on \(X\).  
Fix a value \(x\) of \(X\); then
\[
Z = X + Y = x + Y,
\]
so given \(X=x\), the random variable \(Z\) is just \(Y\) shifted by \(x\).  
For a linear transformation \(U = Y + b\), the PDF satisfies
\[
f_U(u) = f_Y(u-b),
\]
because shifting a continuous random variable by a constant simply shifts its density.  
Thus, for each fixed \(x\),
\[
f_{Z\mid X}(z\mid x)
=
f_Y(z-x).
\]

Now use the multiplication rule for densities to get the joint PDF of \((X,Z)\).  
We have
\[
f_{X,Z}(x,z)
=
f_X(x)\,f_{Z\mid X}(z\mid x)
=
f_X(x)\,f_Y(z-x).
\]
To obtain the marginal density of \(Z\), integrate out \(x\) from the joint density:
\[
f_Z(z)
=
\int_{-\infty}^{\infty} f_{X,Z}(x,z)\,dx
=
\int_{-\infty}^{\infty} f_X(x)\,f_Y(z-x)\,dx.
\]
This is exactly the continuous convolution formula.

\bigskip
\textsc{Graphical interpretation}

The mechanics of using this formula mirror those of the discrete case.  
To visualize the integral, one may take the graph of \(f_Y(y)\), reflect it horizontally to obtain \(f_Y(-y)\), and then shift it by an amount \(z\) to form the function \(y \mapsto f_Y(z-y)\).  
For a fixed \(z\), the integrand \(f_X(x)\,f_Y(z-x)\) multiplies the height of the \(f_X\) graph at \(x\) by the height of the shifted-reflected \(f_Y\) graph at the same \(x\).  
Integrating over all \(x\) adds up these products and yields \(f_Z(z)\).


\subsubsection{Sum of independent normal random variables}

Let \(X\) and \(Y\) be independent normal random variables.  
Assume
\[
X \sim \mathcal{N}(\mu_X,\sigma_X^{2}),
\qquad
Y \sim \mathcal{N}(\mu_Y,\sigma_Y^{2}).
\]
We define
\[
Z = X + Y
\]
and we want to find the PDF and distribution of \(Z\).

\bigskip
\textsc{Using the convolution formula}

Because \(X\) and \(Y\) are independent and continuous, the PDF of \(Z\) is given by the convolution
\[
f_Z(z)
=
\int_{-\infty}^{\infty} f_X(x)\,f_Y(z-x)\,dx.
\]
For normal \(X\) and \(Y\),
\[
f_X(x)
=
\frac{1}{\sqrt{2\pi}\,\sigma_X}
\exp\!\left(-\,\frac{(x-\mu_X)^{2}}{2\sigma_X^{2}}\right),
\]
\[
f_Y(y)
=
\frac{1}{\sqrt{2\pi}\,\sigma_Y}
\exp\!\left(-\,\frac{(y-\mu_Y)^{2}}{2\sigma_Y^{2}}\right).
\]
In the convolution, \(y\) is replaced by \(z-x\), so
\[
f_Z(z)
=
\int_{-\infty}^{\infty}
\frac{1}{\sqrt{2\pi}\,\sigma_X}
\exp\!\left(-\,\frac{(x-\mu_X)^{2}}{2\sigma_X^{2}}\right)
\cdot
\frac{1}{\sqrt{2\pi}\,\sigma_Y}
\exp\!\left(-\,\frac{(z-x-\mu_Y)^{2}}{2\sigma_Y^{2}}\right)
\,dx.
\]
The integrand is a product of two exponentials in \(x\).  
By expanding the squares in the exponents and collecting terms in \(x\), one can rewrite the product as a single Gaussian in \(x\) times a factor that depends only on \(z\).  
Integrating the Gaussian in \(x\) over the real line yields \(1\), leaving a remaining factor in \(z\) that has the form of a normal PDF.

\newpage
\textsc{Resulting distribution}

The final result is that \(Z\) is again normal with
\[
Z \sim \mathcal{N}(\mu_X + \mu_Y,\ \sigma_X^{2} + \sigma_Y^{2}).
\]
Equivalently, the PDF of \(Z\) can be written as
\[
f_Z(z)
=
\frac{1}{\sqrt{2\pi(\sigma_X^{2}+\sigma_Y^{2})}}
\exp\!\left(
-\,
\frac{\big(z-(\mu_X+\mu_Y)\big)^{2}}{2(\sigma_X^{2}+\sigma_Y^{2})}
\right).
\]
The mean of the sum is the sum of the means and the variance of the sum is the sum of the variances.

\bigskip
\textsc{Extension to more than two variables}

Now consider independent normal random variables
\[
X_1, X_2, \dots, X_n,
\qquad
X_i \sim \mathcal{N}(\mu_i,\sigma_i^{2}).
\]
From the two-variable result, the sum \(X_1 + X_2\) is normal.  
Because \(X_3\) is independent of \((X_1 + X_2)\) and normal, the sum \((X_1 + X_2) + X_3\) is also normal.  
Continuing inductively, any finite sum
\[
S_n = X_1 + X_2 + \cdots + X_n
\]
is normal with
\[
S_n \sim \mathcal{N}\!\left(\sum_{i=1}^{n} \mu_i,\ \sum_{i=1}^{n} \sigma_i^{2}\right).
\]
Thus the class of normal distributions is closed under finite sums of independent variables, which makes normal models particularly convenient when studying linear combinations of random signals or errors.


\subsubsection{Covariance and dependence}

Consider two random variables \(X\) and \(Y\) with means \(0\).  
If \(X\) and \(Y\) are independent, then
\[
\mathbb{E}[XY]
=
\mathbb{E}[X]\mathbb{E}[Y]
=
0.
\]
Now imagine a discrete joint distribution where points with \(X>0\) tend to occur together with \(Y>0\), and points with \(X<0\) tend to occur together with \(Y<0\).  
Most outcomes then fall in the quadrants where the product \(XY\) is positive, so the average value \(\mathbb{E}[XY]\) is positive.  
In contrast, if typical outcomes lie in quadrants where \(X\) and \(Y\) have opposite signs, then \(XY\) is usually negative and \(\mathbb{E}[XY]\) becomes negative.  
Thus, in the zero-mean case, the sign of \(\mathbb{E}[XY]\) tells us whether \(X\) and \(Y\) tend to move in the same direction or in opposite directions.

\bigskip
\textsc{General definition}

For general random variables \(X\) and \(Y\) (discrete or continuous) with possibly nonzero means, the covariance is defined as
\[
\operatorname{Cov}(X,Y)
=
\mathbb{E}\big[(X-\mathbb{E}[X])\,(Y-\mathbb{E}[Y])\big].
\]
Here \(X-\mathbb{E}[X]\) and \(Y-\mathbb{E}[Y]\) are the deviations from their respective mean values.  
If the covariance is positive, then when \(X\) is typically above its mean, \(Y\) also tends to be above its mean, and when \(X\) is below its mean, \(Y\) tends to be below its mean.  
If the covariance is negative, then \(X\) being above its mean is usually accompanied by \(Y\) being below its mean, and vice versa.  
In summary, the covariance measures whether two random variables tend to be simultaneously high or low in an average sense.

\bigskip
\textsc{Independence and zero covariance}

If \(X\) and \(Y\) are independent, then
\[
\operatorname{Cov}(X,Y)
=
\mathbb{E}\big[(X-\mathbb{E}[X])\,(Y-\mathbb{E}[Y])\big]
=
\mathbb{E}[X-\mathbb{E}[X]]\;\mathbb{E}[Y-\mathbb{E}[Y]]
=
0,
\]
because \(X-\mathbb{E}[X]\) and \(Y-\mathbb{E}[Y]\) are independent and each has expectation \(0\).  
Thus independence implies zero covariance.  
However, the converse is not true: there exist dependent random variables whose covariance is zero.

\bigskip
\textsc{Example of dependence with zero covariance}

Consider a discrete example with four possible outcomes arranged so that at every outcome either \(X=0\) or \(Y=0\).  
In this construction, the product \(XY\) is always equal to \(0\), so
\[
\mathbb{E}[XY] = 0.
\]
By symmetry, the means of \(X\) and \(Y\) are also \(0\), so
\[
\operatorname{Cov}(X,Y)
=
\mathbb{E}[XY] - \mathbb{E}[X]\mathbb{E}[Y]
=
0.
\]
Nevertheless, \(X\) and \(Y\) are not independent.  
For instance, if you are told that \(X=1\), you immediately know which outcome occurred and can deduce that \(Y=0\) with certainty, so knowledge of \(X\) reveals information about \(Y\).


\subsubsection{Covariance: basic algebraic properties}

For any random variable \(X\), the covariance of \(X\) with itself is
\[
\operatorname{Cov}(X,X)
=
\mathbb{E}\big[(X-\mathbb{E}[X])^{2}\big]
=
\operatorname{Var}(X).
\]
Thus variance is a special case of covariance.

\bigskip
\textsc{Alternative formula}

Starting from the definition
\[
\operatorname{Cov}(X,Y)
=
\mathbb{E}\big[(X-\mathbb{E}[X])\,(Y-\mathbb{E}[Y])\big]
\]
and expanding the product gives
\[
\operatorname{Cov}(X,Y)
=
\mathbb{E}[XY]
-
\mathbb{E}[X]\,\mathbb{E}[Y]
-
\mathbb{E}[X]\,\mathbb{E}[Y]
+
\mathbb{E}[X]\,\mathbb{E}[Y].
\]
Using linearity of expectation and pulling constants outside, the middle two terms cancel, leaving
\[
\operatorname{Cov}(X,Y)
=
\mathbb{E}[XY]
-
\mathbb{E}[X]\,\mathbb{E}[Y].
\]
This identity often simplifies covariance calculations.

\bigskip
\textsc{Effect of linear transformations}

Consider a linear function of \(X\),
\[
U = aX + b
\]
and another random variable \(Y\).  
Using the definition of covariance and linearity of expectation, one can show
\[
\operatorname{Cov}(aX+b,\;Y)
=
a\,\operatorname{Cov}(X,Y).
\]
Multiplying \(X\) by \(a\) scales the covariance by \(a\), while adding a constant \(b\) does not change the covariance.

\bigskip
\textsc{Linearity in one argument}

For random variables \(X\), \(Y\) and \(Z\) (assume for simplicity that they have finite second moments),
\[
\operatorname{Cov}\big(X,\;Y+Z\big)
=
\mathbb{E}\big[X(Y+Z)\big]
-
\mathbb{E}[X]\,\mathbb{E}[Y+Z].
\]
Expanding both terms and collecting gives
\[
\operatorname{Cov}\big(X,\;Y+Z\big)
=
\operatorname{Cov}(X,Y)
+
\operatorname{Cov}(X,Z).
\]
So covariance is linear in each argument when the other is held fixed.



\subsubsection{Variance of a sum of random variables}

Let \(X_1\) and \(X_2\) be random variables and consider their sum
\[
S = X_1 + X_2.
\]
By definition,
\[
\operatorname{Var}(S)
=
\mathbb{E}\Big[\big(S-\mathbb{E}[S]\big)^{2}\Big]
=
\mathbb{E}\Big[\big((X_1+X_2)-(\mathbb{E}[X_1]+\mathbb{E}[X_2])\big)^{2}\Big].
\]
Rewriting inside the square,
\[
S-\mathbb{E}[S]
=
(X_1-\mathbb{E}[X_1]) + (X_2-\mathbb{E}[X_2]).
\]
Expanding the square gives
\[
\big(S-\mathbb{E}[S]\big)^{2}
=
(X_1-\mathbb{E}[X_1])^{2}
+
(X_2-\mathbb{E}[X_2])^{2}
+
2\,(X_1-\mathbb{E}[X_1])(X_2-\mathbb{E}[X_2]).
\]
Taking expectations and using linearity,
\[
\operatorname{Var}(S)
=
\operatorname{Var}(X_1)
+
\operatorname{Var}(X_2)
+
2\,\operatorname{Cov}(X_1,X_2).
\]
If \(X_1\) and \(X_2\) are independent, then \(\operatorname{Cov}(X_1,X_2)=0\) and
\[
\operatorname{Var}(X_1+X_2)
=
\operatorname{Var}(X_1)+\operatorname{Var}(X_2).
\]

\bigskip
\textsc{Many random variables}

Now consider \(n\) random variables \(X_1,\dots,X_n\) and their sum
\[
S_n = X_1 + X_2 + \cdots + X_n.
\]
For simplicity assume \(\mathbb{E}[X_i]=0\), so
\[
\operatorname{Var}(S_n)
=
\mathbb{E}\big[S_n^{2}\big]
=
\mathbb{E}\Big[\Big(\sum_{i=1}^{n} X_i\Big)^{2}\Big].
\]
Expanding the square,
\[
S_n^{2}
=
\sum_{i=1}^{n} X_i^{2}
+
\sum_{\substack{i,j=1\\ i\ne j}}^{n} X_iX_j.
\]
Taking expectations,
\[
\operatorname{Var}(S_n)
=
\sum_{i=1}^{n} \mathbb{E}[X_i^{2}]
+
\sum_{\substack{i,j=1\\ i\ne j}}^{n} \mathbb{E}[X_iX_j].
\]
Because \(\mathbb{E}[X_i]=0\), each \(\mathbb{E}[X_i^{2}] = \operatorname{Var}(X_i)\) and each \(\mathbb{E}[X_iX_j] = \operatorname{Cov}(X_i,X_j)\).  
Thus
\[
\operatorname{Var}(S_n)
=
\sum_{i=1}^{n} \operatorname{Var}(X_i)
+
\sum_{\substack{i,j=1\\ i\ne j}}^{n} \operatorname{Cov}(X_i,X_j).
\]
Equivalently, one can write
\[
\operatorname{Var}(S_n)
=
\sum_{i=1}^{n} \sum_{j=1}^{n} \operatorname{Cov}(X_i,X_j),
\]
since the diagonal terms \(\operatorname{Cov}(X_i,X_i)\) are exactly the variances.  
If all covariances are zero, then the variance of the sum reduces to the sum of the individual variances.  
This happens in particular when the \(X_i\) are independent, but the formula itself holds regardless of whether they are independent or have nonzero covariances.



\subsubsection{The correlation coefficient}

The covariance \(\operatorname{Cov}(X,Y)\) indicates whether two random variables tend to move together, but its numerical value is hard to interpret and it depends on the measurement units.  
If \(X\) and \(Y\) are measured in meters, then the covariance has units of \(\text{meters}^{2}\), which makes comparisons across problems or unit systems difficult.  
To obtain a dimensionless and more interpretable measure of linear dependence, we normalize the covariance by the standard deviations of the two variables.

\bigskip
\textsc{Definition and basic properties}

For random variables \(X\) and \(Y\) with standard deviations \(\sigma_X>0\) and \(\sigma_Y>0\), the correlation coefficient is defined as
\[
\rho_{X,Y}
=
\frac{\operatorname{Cov}(X,Y)}{\sigma_X\,\sigma_Y}.
\]
Because the numerator has the same units as the product \(\sigma_X\sigma_Y\), the ratio \(\rho_{X,Y}\) is dimensionless.  
If either standard deviation is zero, then the corresponding variable is almost surely constant and the correlation coefficient is not defined because the denominator would be zero.  
A fundamental fact is that
\[
-1 \le \rho_{X,Y} \le 1,
\]
which provides an absolute scale for judging the strength of linear association between \(X\) and \(Y\).

\bigskip
\textsc{Interpretation and extreme cases}

If \(X\) and \(Y\) are independent and have finite variances, then \(\operatorname{Cov}(X,Y)=0\) and
\[
\rho_{X,Y}=0,
\]
and in this case the variables are said to be \textit{uncorrelated}.  
The converse is not generally true: there exist dependent random variables with zero covariance and therefore zero correlation.  
At the other extreme, consider the correlation coefficient of a random variable with itself:
\[
\rho_{X,X}
=
\frac{\operatorname{Cov}(X,X)}{\sigma_X^{2}}
=
\frac{\operatorname{Var}(X)}{\sigma_X^{2}}
=
1.
\]
Similarly, the correlation coefficient of \(X\) with \(-X\) is \(-1\).  
More generally, if there is a deterministic linear relationship between \(X\) and \(Y\) of the form
\[
Y = aX + b
\]
with probability one, then
\[
|\rho_{X,Y}| = 1,
\]
with \(\rho_{X,Y}=1\) when \(a>0\) and \(\rho_{X,Y}=-1\) when \(a<0\).  
Conversely, if \(|\rho_{X,Y}|=1\), then \(X\) and \(Y\) are related by such a deterministic linear formula, so extreme correlations \(\pm1\) correspond exactly to perfect linear dependence.

\bigskip
\textsc{Effect of linear transformations}

Let \(U = aX + b\) for real constants \(a\) and \(b\) and consider the correlation between \(U\) and \(Y\).  
From covariance properties,
\[
\operatorname{Cov}(aX+b,\;Y)
=
a\,\operatorname{Cov}(X,Y).
\]
The standard deviation of \(aX+b\) is \(|a|\,\sigma_X\), because shifting by \(b\) does not change the spread and scaling by \(a\) multiplies the standard deviation by \(|a|\).  
Therefore,
\[
\rho_{aX+b,\;Y}
=
\frac{\operatorname{Cov}(aX+b,\;Y)}{\sigma_{aX+b}\,\sigma_Y}
=
\frac{a\,\operatorname{Cov}(X,Y)}{|a|\sigma_X\,\sigma_Y}
=
\operatorname{sign}(a)\,\rho_{X,Y},
\]
where
\[
\operatorname{sign}(a)
=
\begin{cases}
1, & a>0,\\[4pt]
-1, & a<0.
\end{cases}
\]
The magnitude \(|\rho_{aX+b,\;Y}|\) is the same as \(|\rho_{X,Y}|\); only the sign may change if \(a<0\).  
Thus changing units for a variable (for example, converting a temperature from degrees Celsius to degrees Fahrenheit) does not change the absolute value of the correlation with another variable, reflecting the fact that correlation is unitless and invariant under positive linear rescalings.


\subsubsection{Key properties of the correlation coefficient}

\bigskip
\textsc{Bounded between \(-1\) and \(1\)}

To justify the basic bound on the correlation coefficient, consider a simplified setting in which
\[
\mathbb{E}[X]=\mathbb{E}[Y]=0,
\qquad
\operatorname{Var}(X)=\operatorname{Var}(Y)=1.
\]
In this case the standard deviations are \(1\), so the correlation coefficient reduces to
\[
\rho = \rho_{X,Y} = \mathbb{E}[XY].
\]
Define the random variable
\[
W = X - \rho Y
\]
and examine the expectation of its square:
\[
\mathbb{E}[W^{2}]
=
\mathbb{E}\big[(X-\rho Y)^{2}\big].
\]
Expanding the square gives
\[
\mathbb{E}[W^{2}]
=
\mathbb{E}[X^{2}]
-
2\rho\,\mathbb{E}[XY]
+
\rho^{2}\,\mathbb{E}[Y^{2}].
\]
Under our assumptions \(\mathbb{E}[X^{2}] = \operatorname{Var}(X) = 1\) and \(\mathbb{E}[Y^{2}] = 1\).  
Also \(\mathbb{E}[XY]=\rho\).  
Substituting,
\[
\mathbb{E}[W^{2}]
=
1
-
2\rho^{2}
+
\rho^{2}
=
1-\rho^{2}.
\]
Since \(W^{2}\) is always nonnegative, its expectation must be nonnegative:
\[
\mathbb{E}[W^{2}] \ge 0
\quad\Longrightarrow\quad
1-\rho^{2} \ge 0.
\]
Therefore \(\rho^{2} \le 1\), which implies
\[
-1 \le \rho \le 1.
\]
With a bit more algebra the same conclusion holds for general random variables \(X\) and \(Y\) with arbitrary means and nonzero variances.

\bigskip
\textsc{Characterizing the extreme case \(|\rho|=1\)}

Suppose now that the absolute value of the correlation coefficient is equal to \(1\).  
In the setting above, this means \(\rho^{2}=1\), so from the previous result
\[
\mathbb{E}[W^{2}]
=
1-\rho^{2}
=
0.
\]
Again \(W^{2}\) is nonnegative, and a nonnegative random variable with zero expectation must be zero with probability \(1\).  
Hence
\[
W^{2} = 0 \text{ with probability }1,
\quad\text{so}\quad
X-\rho Y = 0 \text{ with probability }1.
\]
This shows that
\[
X = \rho Y
\quad\text{with probability }1,
\]
which is a deterministic linear relation between \(X\) and \(Y\).  
Because \(|\rho|=1\), there are two possibilities:
\[
\rho = 1 \quad\Longrightarrow\quad X = Y \text{ almost surely},
\]
\[
\rho = -1 \quad\Longrightarrow\quad X = -Y \text{ almost surely}.
\]
Thus \(|\rho|=1\) corresponds to an extreme form of dependence, in which one random variable is an exact linear function of the other.


\subsubsection{Interpreting the correlation coefficient}

A correlation coefficient measures the strength of a \textit{linear association} between two random variables, not a causal relationship.  
A value such as \(\rho_{X,Y} = 0.5\) indicates that high values of \(X\) tend to be accompanied by high values of \(Y\) and low values of \(X\) tend to be accompanied by low values of \(Y\), in an average sense.  
However, this does not mean that changing \(X\) will directly cause a change in \(Y\), or vice versa.  
In many situations a nonzero correlation arises because both variables are influenced by some common underlying factor, rather than one variable directly affecting the other.  
For example, mathematical aptitude and musical skill may be positively correlated because both are influenced by certain shared cognitive traits, even though studying one does not necessarily cause improvement in the other.

\bigskip
\textsc{A simple model with a hidden common factor}

To make this idea concrete, consider independent random variables \(Z\), \(V\), and \(W\) with
\[
\mathbb{E}[Z] = \mathbb{E}[V] = \mathbb{E}[W] = 0,
\qquad
\operatorname{Var}(Z) = \operatorname{Var}(V) = \operatorname{Var}(W) = 1.
\]
Define
\[
X = Z + V,
\qquad
Y = Z + W.
\]
Here the underlying factor \(Z\) affects both \(X\) and \(Y\), while \(V\) and \(W\) are independent “noise” terms specific to each variable.  
There is no direct causal influence from \(X\) to \(Y\) or from \(Y\) to \(X\); the dependence comes entirely from the shared component \(Z\).

\bigskip
\textsc{Computing the correlation in the example}

First compute the variances.  
Because \(Z\) and \(V\) are independent,
\[
\operatorname{Var}(X)
=
\operatorname{Var}(Z) + \operatorname{Var}(V)
=
1 + 1
=
2,
\]
and similarly
\[
\operatorname{Var}(Y)
=
\operatorname{Var}(Z) + \operatorname{Var}(W)
=
2.
\]
Thus the standard deviations are
\[
\sigma_X = \sigma_Y = \sqrt{2}.
\]
Next compute the covariance.  
Since \(\mathbb{E}[X]=\mathbb{E}[Y]=0\),
\[
\operatorname{Cov}(X,Y)
=
\mathbb{E}[XY]
=
\mathbb{E}[(Z+V)(Z+W)]
=
\mathbb{E}[Z^{2}]
+
\mathbb{E}[ZW]
+
\mathbb{E}[VZ]
+
\mathbb{E}[VW].
\]
Because \(Z\), \(V\), and \(W\) are independent and have zero means,
\[
\mathbb{E}[Z^{2}] = 1,
\qquad
\mathbb{E}[ZW] = \mathbb{E}[Z]\mathbb{E}[W] = 0,
\qquad
\mathbb{E}[VZ] = 0,
\qquad
\mathbb{E}[VW] = 0,
\]
so
\[
\operatorname{Cov}(X,Y) = 1.
\]
The correlation coefficient is then
\[
\rho_{X,Y}
=
\frac{\operatorname{Cov}(X,Y)}{\sigma_X\sigma_Y}
=
\frac{1}{\sqrt{2}\,\sqrt{2}}
=
\frac{1}{2}.
\]

\bigskip
\textsc{Qualitative meaning of \(\rho = \tfrac{1}{2}\)}

In this model, the common factor \(Z\) and the idiosyncratic parts \(V\) and \(W\) all have the same variance, so the “shared” part and the “individual noise” parts have comparable strength.  
The resulting correlation \(\rho=\tfrac{1}{2}\) reflects a moderate linear association: \(X\) and \(Y\) tend to move together because of their common component \(Z\), but not perfectly, since the independent noises \(V\) and \(W\) partly obscure that relation.  
If the noise terms \(V\) and \(W\) were absent, we would have \(X=Y=Z\) and the correlation would be \(1\).  
If instead \(V\) and \(W\) had very large variances compared to \(Z\), their idiosyncratic fluctuations would dominate, the influence of \(Z\) would be hard to detect, and the correlation would be much smaller, approaching \(0\) in the limit.


\subsubsection{Why correlations matter in practice}

Consider an investment company with capital of \(100\) million that is allocated equally across \(10\) different states.  
Let \(X_i\) denote the random profit (in millions) from state \(i\), so each investment is \(10\) million and
\[
\mathbb{E}[X_i] = 1,
\qquad
\operatorname{SD}(X_i) = 1.3,
\qquad
\operatorname{Var}(X_i) = 1.3^{2} = 1.69.
\]
The total profit is
\[
S = X_1 + X_2 + \cdots + X_{10}.
\]
Its expectation is
\[
\mathbb{E}[S] = \sum_{i=1}^{10} \mathbb{E}[X_i] = 10.
\]
If the state markets are truly unrelated so that the \(X_i\) are uncorrelated, then all covariances are zero and
\[
\operatorname{Var}(S)
=
\sum_{i=1}^{10} \operatorname{Var}(X_i)
=
10 \times 1.69
=
16.9.
\]
Thus
\[
\operatorname{SD}(S) = \sqrt{16.9} \approx 4.1.
\]
Since the mean profit is \(10\), breaking even corresponds to an outcome \(2.5\) standard deviations below the mean (because \(10 - 2.5 \times 4.1 \approx 0\)).  
Outcomes this far below the mean have relatively small probability for many light-tailed distributions, so in the uncorrelated case diversification gives a high chance of a positive total return.

\bigskip
\textsc{Effect of strong positive correlations}

Suppose instead that the state markets are strongly linked so that
\[
\rho_{X_i,X_j} = 0.9
\quad\text{for all } i\ne j.
\]
Each covariance is then
\[
\operatorname{Cov}(X_i,X_j)
=
\rho_{X_i,X_j}\,\operatorname{SD}(X_i)\,\operatorname{SD}(X_j)
=
0.9 \times 1.3 \times 1.3
=
1.52.
\]
The variance of the sum becomes
\[
\operatorname{Var}(S)
=
\sum_{i=1}^{10}\operatorname{Var}(X_i)
+
\sum_{\substack{i,j=1\\ i\ne j}}^{10}\operatorname{Cov}(X_i,X_j)
=
10 \times 1.69
+
90 \times 1.52.
\]
Here the first term accounts for the \(10\) individual variances, and the second term has \(90\) covariance entries because there are \(10\cdot 9\) ordered pairs \((i,j)\) with \(i\ne j\).  
Carrying out the arithmetic,
\[
\operatorname{Var}(S)
=
16.9 + 136.8
=
153.7 \approx 154,
\]
so
\[
\operatorname{SD}(S)
\approx
\sqrt{154}
\approx
12.4.
\]
The expected profit is still \(10\), but now a one-standard-deviation drop from the mean (\(10-12.4\)) already leads to a negative total return.  
Such a deviation has substantial probability, so the overall investment becomes much riskier when returns are highly positively correlated.

\bigskip
\textsc{Lessons for risk and diversification}

This example illustrates that diversification reduces risk only when the underlying positions are not strongly positively correlated.  
If the component returns move largely together because of common economic forces, the covariance terms dominate the variance of the aggregate portfolio and the overall risk can remain large even when the portfolio includes many different assets.  
In the context of real estate, it is not enough to invest in multiple regions if housing markets across regions are driven by shared macroeconomic factors that create high correlations.  
When such correlations are underestimated or ignored, as in parts of the financial sector before the global financial crisis, portfolios that appear well diversified can still suffer large simultaneous losses.